% concatenated from main.tex
\documentclass{amsart}
\usepackage[margin=1in]{geometry}
\usepackage[square,numbers]{natbib}
\usepackage{shuffle}
\usepackage{mathrsfs}
\usepackage{bbm}
\usepackage{ifthen}
\usepackage{tikz}
\usetikzlibrary{matrix,arrows,calc}
\usepackage{mathdots}
\usepackage{array}
\usepackage{multirow}
\usepackage{cjhebrew}
\usepackage{hyperref}
\usepackage{float}
\usepackage{caption}
\usepackage{amscd}
\usepackage{tabularx}
\usepackage{extarrows}
\usepackage{ytableau}
\usepackage{amsmath}
\usepackage{amsthm}
\usepackage{amsfonts}
\usepackage{amssymb}
\usepackage{mathtools}
\usepackage{soul}
\usepackage{color}
\newcommand{\hlc}[2][yellow]{ {\sethlcolor{#1} \hl{#2}} }

\newtheorem{theorem}{Theorem}[section]
\newtheorem{conjecture}[theorem]{Conjecture}
\newtheorem{proposition}[theorem]{Proposition}
\newtheorem{corollary}[theorem]{Corollary}
\newtheorem{lemma}[theorem]{Lemma}
\theoremstyle{definition}
\newtheorem{definition}[theorem]{Definition}
\newtheorem{example}[theorem]{Example}
\newtheorem{remark}[theorem]{Remark}
\newtheorem{problem}[theorem]{Problem}

\def\unprotectedboldentry#1{\textcolor{Red}{\textbf{#1}}} 
\def\boldentry{\protect\unprotectedboldentry} 

\newcommand{\tikztableau}[2][scale=0.6,every node/.style={font=\small}]{
 \def\newtableau{#2} 
 \begin{array}{c} 
 \begin{tikzpicture}[#1] 
 \coordinate (x) at (-0.5,0.5); 
 \coordinate (y) at (-0.5,0.5); 
 \foreach \row in \newtableau { 
 \coordinate (x) at ($(x)-(0,1)$); 
 \coordinate (y) at (x); 
 \foreach \entry in \row { 
 \ifthenelse{\equal{\entry}{X}} 
 { 
 \node (y) at ($(y) + (1,0)$) {}; 
 \fill[color=lightgray] ($(y)-(0.5,0.5)$) rectangle +(1,1); 
 \draw[color=black] ($(y)-(0.5,0.5)$) rectangle +(1,1); 
 } 
 { 
 \ifthenelse{\equal{\entry}{\boldentry X}} 
 { 
 \node (y) at ($(y) + (1,0)$) {}; 
 \fill[color=darkgray] ($(y)-(0.5,0.5)$) rectangle +(1,1); 
 \draw ($(y)-(0.5,0.5)$) rectangle +(1,1); 
 } 
 { 
 \node (y) at ($(y) + (1,0)$) {\entry}; 
 \draw ($(y)-(0.5,0.5)$) rectangle +(1,1);} } 
 }} 
 \end{tikzpicture} 
 \end{array}} 
 
\newcommand{\tikztableausmall}[1]{\tikztableau[scale=0.45,every node/.style={font=\rm\small}]{#1}}
\newcommand{\doi}[1]{\href{http://dx.doi.org/#1}{\texttt{doi:#1}}}
\newcommand{\arxiv}[1]{\href{http://arxiv.org/abs/#1}{\texttt{arXiv:#1}}}

\newcommand{\C}{\mathbb{C}}
\newcommand{\N}{\mathbb{N}}
\newcommand{\Z}{\mathbb{Z}}
\newcommand{\Max}{Max_{\succeq}}
\newcommand{\Q}{\mathbb{Q}}
\newcommand{\R}{\mathbb{R}}
\newcommand{\tr}[1]{\textcolor{red}{ #1}}
\newcommand{\lx}{\mathfrak{L}}
\newcommand{\os}[2]{IO\Pi_{#1, #2}}
\newcommand{\co}{\overrightarrow{Co}}
\newcommand{\ts}{\overrightarrow{TS}}

%\def\shin{{\text{\large\</s>}}}

\title{Lexical tableaux and quasisymmetric functions}

\author[J. M. Campbell]{John M. Campbell}\address{Department of Mathematics and Statistics, Dalhousie University, Halifax, NS}\email{jh241966@dal.ca}

\author[S. Daugherty]{Spencer Daugherty}\address{Department of Mathematics, University of Colorado Boulder, Boulder, CO}\email{spencer.daugherty@colorado.edu}

\begin{document}

\maketitle

\begin{abstract}
 There is a natural bijection between standard immaculate tableaux of composition shape $\alpha \vDash n$ and length $\ell(\alpha) = 
 k$ and the $ \left\{ \begin{smallmatrix} n \\ k \end{smallmatrix} \right\} $ set-partitions of $\{ 1, 2, \ldots, n \}$ into $k$ blocks, for the 
 Stirling number $ \left\{ \begin{smallmatrix} n \\ k \end{smallmatrix} \right\} $ of the second kind. We introduce a family of tableaux that 
 we refer to as \emph{lexical tableaux} that generalize immaculate tableaux in such a way that there is a bijection between standard 
 lexical tableaux of shape $\alpha \vDash n$ and length $\ell(\alpha) = k$ and the $ \left[ \begin{smallmatrix} n \\ k \end{smallmatrix} 
 \right] $ permutations on $\{ 1, 2, \ldots, n \}$ with $k$ disjoint cycles. In addition to the entries in the first column strictly increasing, the defining characteristic of lexical tableaux is that the word $w$ formed by the consecutive labels in any row is 
 the lexicographically smallest out of all cyclic permutations of $w$. This includes weakly increasing words, and thus  lexical tableaux provide a natural generalization of immaculate tableaux. Extending this generalization, we introduce a pair of dual bases of the Hopf algebras 
 $\textsf{QSym}$ and $\textsf{NSym}$ defined in terms of lexical tableaux. We present two expansions of these bases, involving the monomial and fundamental bases (or,  dually,  the ribbon and complete homogeneous bases), using Kostka coefficient analogues and coefficients derived from standard lexical tableaux. 
\end{abstract}

\vspace{0.1in}

\noindent {\footnotesize{\emph{MSC:} 05E05, 05E10}}

\vspace{0.1in}

\noindent {\footnotesize{\emph{Keywords:} tableau, quasisymmetric function, symmetric function, noncommutative symmetric function, symmetric group, tabloid, permutation, integer partition, integer composition, Hopf algebra, Young tableau, immaculate tableau}}

\section{Introduction}
 An \emph{integer partition} is a weakly decreasing and finite tuple $\lambda = (\lambda_1, \lambda_2, \ldots, \lambda_{\ell(\lambda)})$ 
 of positive integers, where $|\lambda|$ denotes the sum of the parts of $\lambda$, and we write $\lambda \vdash n$ if $|\lambda| = 
 n$. Similarly, an \emph{integer composition} $\alpha = (\alpha_1, \alpha_2, \ldots, \alpha_{\ell(\alpha)})$ is a finite tuple of positive 
 integers, where we write $|\alpha|$ in place of the   sum of the entries of $\alpha$, and we write $\alpha \vDash n$ if $|\alpha| = n$. A  
 \emph{tableau}, for the purposes of this paper,   may be understood as a two-dimensional arrangement of (labeled or unlabelled) cells 
 that are positioned into left-justified rows,   with the \emph{shape} of a tableau $T$ forming an integer composition $\alpha$ such that  
 the number of cells in the $i^{\text{th}}$   row of $T$ (from the bottom according to the so-called French convention) is equal to  
 $\alpha_i$ for $i \in \{ 1, 2, \ldots, \ell(\alpha) \}$.   Tableaux and tableaux-like objects are of fundamental importance in the representation  
 theory of the symmetric group and many   related areas of algebraic combinatorics. \emph{Young tableaux}, in particular, are especially  
 significant in the study of and application of   symmetric group representations. A family of Young-like tableaux that arose in the  
 construction of a noncommutative analogue of the   Schur symmetric functions are the \emph{immaculate tableaux} introduced in a  
 seminal paper by Berg et  al.~\cite{BergBergeronSaliolaSerranoZabrocki2014}.  Our explorations based on the combinatorics of  
 immaculate tableaux have led us to   generalize such tableaux via a family of tableaux that we refer to as \emph{lexical tableaux} and that  
 we apply to introduce new bases of   the Hopf algebras $\textsf{QSym}$ and $\textsf{NSym}$ (reviewed in Section  
 \ref{sectionPreliminaries} below) of quasisymmetric  functions and of noncommutative symmetric functions. 

 One of the most significant results in the representation theory of the symmetric group is that the isomorphism classes of the simple 
 $\mathbb{C}S_n$-modules are in bijection with partitions $\lambda \vdash n$, and, moreover, that the dimension and multiplicity of 
 the irreducible $\mathbb{C}S_n$-module corresponding to $\lambda$ is equal to the number $f^{\lambda}$ of standard Young tableaux 
 of shape $\lambda$. This raises questions as to how similar properties could be obtained with the use of composition tableaux in place 
 of partition tableaux, and immaculate tableaux can be thought of as arising in this way. The interest in the study of combinatorial 
 properties associated with standard immaculate tableaux is evidenced by Gao and Yang's bijective proof of the hook-length formula for 
 standard immaculate tableaux \cite{GaoYang2016} together with Sun and Hu's probabilistic method for determining the number of 
 standard immaculate tableaux of a given shape \cite{SunHu2018}. 

 Since evaluations of finite summations involving $f^{\lambda}$ often arise in the context of applications of Young tableaux, this raises 
 questions as to what would be appropriate as analogues of such evaluations involving standard immaculate tableaux, 
 letting $g^{\alpha}$ denote the number of standard immaculate tableaux 
 (reviewed in Section \ref{sectionPreliminaries} below) of a composition shape $\alpha$.
 In this direction, 
 the summation
\begin{equation}\label{krows}
 \sum_{\substack{ \lambda \vdash n \\ \ell(\lambda)=k}} f^{\lambda} = \text{$\#$ of Young tableaux with $n$ cells and $k$ rows}, 
\end{equation} 
 which gives rise to a number triangle indexed in the Online Encyclopedia of Integer Sequences as {\tt A047884}, has led us to 
 experimentally discover, using the OEIS, a property (given in Theorem \ref{thm:snk} below) 
 concerning a corresponding sum for immaculate tableaux
 given by replacing $f^{\lambda}$ with $g^{\alpha}$ and summing over compositions $\alpha \vDash n$
 of a fixed length $k$. 
 This, in turn, has led us to construct a new family of immaculate-like tableaux.
 
 Given a construction involving all possible permutations of a set of objects, it is natural to consider a corresponding construction whereby 
 the permutations involved are required to be cyclic. For example, the abelian complexity function on infinite words counts subwords up 
 to \emph{all} possible permutations of characters, whereas the cyclic complexity function introduced in 2017 
 \cite{CassaigneFiciSciortinoZamboni2017} counts subwords up to \emph{cyclic} permutations of characters. The way our definition of a 
 lexical tableau, as given in Section \ref{sectionLexical} below, 
 relates to that of an immaculate tableau may be seen by analogy with how the definition of a cyclic complexity function 
 relates to that of an abelian complexity function. Similarly, the way the \emph{cyclic quasisymmetric functions} introduced by Adin et 
 al.\ \cite{AdinGesselReinerRoichman2021} are defined via an invariance property associated with cyclic permutations, relative to the 
 corresponding invariance property for symmetric functions holding for \emph{all} possible permutations, further illustrates how the 
 definition of a lexical tableau provides a natural generalization of immaculate tableaux. Indeed, our construction of lexical tableaux 
 makes use of cyclic shifts by direct analogy with the work of Adin et al.\ \cite{AdinGesselReinerRoichman2021}. As the term \emph{lexical 
 tableau} suggests, there is a close connection between the study of such tableau and the field combinatorics on words. 

 Given a property associated with Schur or immaculate functions, by deriving an analogous identity using lexical tableaux, this, ideally, 
 could help to shed light on the use of new methods that could be applied toward unsolved problems related to the Schur and immaculate 
 bases and representation-theoretic uses of these bases. The problem of extending immaculate and dual immaculate functions using 
 the combinatorial objects involved in the construction of $\{ \mathfrak{S}_{\alpha} \}_{\alpha}$ and its dual is motivated by much in the way 
 of past research on immaculate and dual immaculate functions, including research on the indecomposable modules for the dual 
 immaculate basis \cite{BergBergeronSaliolaSerranoZabrocki2015}, Pieri rules for dual immaculate functions and generalizations 
 \cite{BergeronSanchezOrtegaZabrocki2016,Li2018,NieseSundaramvanWilligenburgWang2024}, multiplicative structures of the immaculate 
 basis \cite{BergBergeronSaliolaSerranoZabrocki2017,Li2018}, the expansion of dual immaculate functions into {Y}oung quasisymmetric 
 {S}chur functions \cite{AllenHallamMason2018}, noncommutative {Bell} polynomials \cite{NovelliThibonToumazet2018}, 
 a generalization of the dual immaculate basis to the polynomial algebra \cite{MasonSearles2021}, the immaculate inverse {K}ostka 
 matrix \cite{LoehrNiese2021}, and a generalization of dual immaculate functions using partially commutative 
 variables \cite{Daugherty2024}.
 
  We introduce  lexical tableaux in
  Section \ref{sectionLexical}, and, in Section \ref{subsectionEnumeration}, we present basic enumerative properties of lexical tableaux.   In Section \ref{sectionqsym}, we define the dual lexical functions in $\textsf{QSym}$ and establish via their monomial expansions that they constitute a basis. This expansion uses an analogue of Kostka coefficients that count certain lexical tableaux. We define the lexical functions in $\textsf{NSym}$ dually. Here, we give a positive expansion of the ribbon noncommutative symmetric functions into the lexical functions. We also present results on the antipodes of lexical basis elements. To close, we present various open problems and directions for future research. 

\section{Preliminaries}\label{sectionPreliminaries}
 We highlight Macdonald's text \cite{Macdonald1995} as the usual monograph on symmetric functions. A review of preliminaries on 
 symmetric functions, as below, is required for our purposes. 

 The rings and algebras considered in this paper will be over $\mathbb{Q}$ for convenience and by convention. Letting the symmetric 
 group $S_n$ act on the polynomial ring $\mathbb{Q}[x_1, x_2, \ldots, x_n]$ by permuting the variables, we let
\begin{equation}\label{Symsuper}
 \textsf{Sym}^{(n)} 
 = \mathbb{Q}[x_1, x_2, 
 \ldots, x_n]^{S_{n}} 
\end{equation} 
 denote the polynomial subring given by polynomials invariant under the action of $S_n$. By letting $\textsf{Sym}^{(n)}_{k}$ consist 
 of the zero polynomial and the homogeneous symmetric polynomials that are of degree $k$, we obtain a graded ring
 structure on \eqref{Symsuper}, with $$ \textsf{Sym}^{(n)} 
 = \bigoplus_{k=0}^{\infty} \textsf{Sym}^{(n)}_{k}.$$ We then take an inverse limit
\begin{equation*} 
 \textsf{Sym}_{k} := \lim_{\substack{\longleftarrow\\ n }} \textsf{Sym}_{k}^{(n)}, 
\end{equation*}
 referring to Macdonald's text for details \cite[p.\ 18]{Macdonald1995}. This allows us to define the \emph{algebra of symmetric 
 functions} $\textsf{Sym}$ so that 
\begin{equation}\label{gradedSym}
 \textsf{Sym} := \bigoplus_{k=0}^{\infty} \textsf{Sym}_{k}.
\end{equation}

 For a given integer partition $\lambda$, we set
\begin{equation*}
 m_{\lambda} = \sum_{ i_{1} < i_{2} < \cdots < i_{\ell(\lambda)} } 
 x_{i_{1}}^{\lambda_{1}} x_{i_{2}}^{\lambda_{2}} \cdots 
 x_{i_{\ell(\lambda)}}^{\lambda_{\ell(\lambda)}}.
\end{equation*}
 Letting $\mathcal{P}$ denote the set of all integer partitions, we may thus define the \emph{monomial basis} of $\textsf{Sym}$ as 
 $\{ m_{\lambda} \}_{\lambda \in \mathcal{P}}$. By then setting $e_0 = 1$ and $ e_{r} = m_{(1^r)} $ for $r > 0$, we then set 
 $e_{\lambda} = e_{\lambda_1} e_{\lambda_2} \cdots e_{\lambda_{\ell(\lambda)}}$ for $\lambda \in \mathcal{P}$, giving us the 
 \emph{elementary basis} $\{ e_{\lambda} \}_{\lambda \in \mathcal{P}}$ of $\textsf{Sym}$. By then setting $h_0 = 1$ and $ h_r = 
 \sum_{\lambda \vdash r} m_{\lambda} $ for $r > 0$, we then set $h_{\lambda} = h_{\lambda_{1}} h_{\lambda_{2}} \cdots 
 h_{\lambda_{\ell(\lambda)}}$ for an integer partition $\lambda$, giving rise to the \emph{complete homogeneous basis} $\{ 
 h_{\lambda} \}_{\lambda \in \mathcal{P}}$ of $\textsf{Sym}$. By then setting $p_r = m_{(r)}$ for $r \geq 1$, we then set $p_{\lambda} = 
 p_{\lambda_{1}} p_{\lambda_{2}} \cdots p_{\lambda_{\ell(\lambda)}}$, and this gives rise to the \emph{power sum basis} $\{ p_{\lambda} 
 \}_{\lambda \in \mathcal{P}}$ of $\textsf{Sym}$. 

 A \emph{semistandard Young tableau} is a tableau of a given partition shape $\lambda$ with each cell labeled with a positive integer 
 and with weakly increasing rows and strictly increasing columns. The \emph{content} of a SSYT $T$ is the finite tuple $t = (t_{1}, t_{2}, 
 \ldots, t_{\ell(t)})$ such that $\ell(t)$ is the maximal label in $T$ and such that the number of labels equal to $i$ is $t_i$ for $i \in \{ 1, 2, 
 \ldots, \ell(t) \}$. We let the \emph{Kostka coefficient} $K_{\lambda, \mu}$ be defined as the number of SSYTs of shape $\lambda$ 
 and content $\mu$. For the case whereby $\mu = (1^{|\lambda|})$, a SSYT of shape $\lambda$ and content $\mu$ is said to be 
 \emph{standard}, and, as above, the number of standard Young tableaux of shape $\lambda$ is denoted as $f^{\lambda}$. 

\begin{example}\label{exampleSY}
 For $\lambda = (2, 2, 1)$, we find that there are $f^{\lambda} = 5$
 standard Young tableaux of the specified shape, as below: 
 $$ \ytableausetup {centertableaux,boxsize=normal,notabloids}
 \begin{ytableau} 3 \\ 2 & 5 \\ 1 & 4 \end{ytableau} \ \ \ \ \ \ \ 
 \ytableausetup {centertableaux,boxsize=normal,notabloids}
 \begin{ytableau} 4 \\ 2 & 5 \\ 1 & 3 \end{ytableau} \ \ \ \ \ \ \ \ytableausetup {centertableaux,boxsize=normal,notabloids}
 \begin{ytableau} 4 \\ 3 & 5 \\ 1 & 2 \end{ytableau} \ \ \ \ \ \ \ \ytableausetup {centertableaux,boxsize=normal,notabloids}
 \begin{ytableau} 5 \\ 2 & 4 \\ 1 & 3 \end{ytableau} \ \ \ \ \ \ \ \ytableausetup {centertableaux,boxsize=normal,notabloids}
 \begin{ytableau} 5 \\ 3 & 4 \\ 1 & 2 \end{ytableau}. $$ 
\end{example}

 This allows us to define the \emph{Schur basis} $\{ s_{\lambda} \}_{\lambda \in \mathcal{P}}$ so that $ h_{\mu} = \sum_{\lambda} 
 K_{\lambda, \mu} s_{\lambda}$. Equivalently, Schur symmetric functions may be defined according to the \emph{Pieri rule} such that
\begin{equation}\label{originalPieri}
 s_{\lambda} h_{r}
 = \sum_{\mu} s_{\mu}, 
\end{equation}
 where the sum in \eqref{originalPieri} is over all partitions $\mu$
 such that the diagram
 for $\mu$ can be obtained
 from that of $\lambda$ by adding $r$ boxes to the diagram of $\lambda$ (so that the added boxes are adjacent to the border of the diagram of $\lambda$ and are otherwise outside of this diagram) and in such a way so that no two boxes are added to the same column. 

 For an integer composition $\alpha$, we write
\begin{equation}\label{quasimonomial}
 M_{\alpha} = \sum_{i_{1} < i_{2} < \cdots < i_{\ell(\alpha)}}
 x_{i_{1}}^{\alpha_{1}} 
 x_{i_{2}}^{\alpha_{2}} \cdots 
x_{i_{\ell(\alpha)}}^{\alpha_{\ell(\alpha)}}. 
\end{equation}
 By then setting $\textsf{QSym}_{k} := \mathscr{L}\{ M_{\alpha} : \alpha \vDash k \}$, we then form a graded algebra, by analogy with 
 \eqref{gradedSym}, by setting 
\begin{equation}\label{gradedQSym}
 \textsf{QSym} :=
 \bigoplus_{k=0}^{\infty} 
 \textsf{QSym}_{k}.
\end{equation}
 The algebra in \eqref{gradedQSym} is referred to as the \emph{algebra of quasisymmetric functions}, and elements in the basis $\{ 
 M_{\alpha} \}_{\alpha \in \mathcal{C}}$ of $\textsf{QSym}$ are referred to as \emph{monomial quasisymmetric functions}. For 
 compositions $\alpha$ and $\beta$, we write $\alpha \succeq \beta$ if $\alpha$ can be obtained by adding together consecutive parts 
 of $\beta$. This allows us to define the \emph{fundamental basis} $\{ F_{\alpha} \}_{\alpha \in \mathcal{C}}$ of $\textsf{QSym}$ so that 
 $ F_{\alpha} = \sum_{\alpha \succeq \beta} M_{\beta}$. Letting $\text{sort}(\alpha)$ denote the integer partition obtained from $\alpha$ 
 by sorting the entries of $\alpha$, we find that $\textsf{Sym}$ is contained in $\textsf{QSym}$, according to the relation $$ m_{\lambda} 
 = \sum_{ \substack{\alpha \in \mathcal{C} \\ \text{sort}(\alpha) = \lambda} } M_{\alpha}. $$ Quasisymmetric functions were introduced by 
 Gessel in 1984 \cite{Gessel1984} and provide deep and major areas of study within algebraic combinatorics. An equivalent version of the 
 algebra $\textsf{NSym}$ dual to $\textsf{QSym}$ was introduced by Gelfand et al.\ in 1995 
 \cite{GelfandKrobLascouxLeclercRetakhThibon1995}, and the importance of $\textsf{NSym}$ within algebraic combinatorics is much like 
 that of its dual $\textsf{QSym}$. Setting $\textsf{NSym}_{k} := \mathscr{L}\{ H_{\alpha} : \alpha \vDash k \}$, letting $H_{\alpha}$ be seen 
 as a variable, we form the graded algebra $$ \textsf{NSym} := \bigoplus_{k=0}^{\infty} \textsf{NSym}_k $$ endowed with the 
 multiplicative operation such that $H_{\alpha} H_{\beta} = H_{\alpha \cdot \beta}$ for the concatenation $\alpha \cdot \beta$ of 
 $\alpha$ and $\beta$. The elementary basis $\{ E_{\alpha} \}_{\alpha \in \mathcal{C}}$ of $\textsf{NSym}$ may then be defined 
 according to the recursion $ E_{n} = \sum_{i=1}^{n} (-1)^{n+1} H_{i} E_{n-i}$, with $E_{\alpha} = E_{\alpha_1} E_{\alpha_2} \cdots 
 E_{\alpha_{\ell(\alpha)}}$, and the ribbon basis $ \{ R_{\alpha} \}_{\alpha \in \mathcal{C}} $ may be defined so that $ R_{\alpha} 
 = \sum_{\beta \succeq \alpha} (-1)^{\ell(\alpha) - \ell(\beta)} H_{\beta}$. 

 The duality between $\textsf{NSym}$ and $\textsf{QSym}$ may be demonstrated using the bases defined above, according to the 
 bilinear pairing $\langle \cdot, \cdot \rangle\colon \textsf{NSym} \times \textsf{QSym} \to \mathbb{Q}$ such that $\langle H_{\alpha}, 
 M_{\beta} \rangle = \delta_{\alpha, \beta}$ for the Kronecker delta function $\delta_{\cdot, \cdot}$, or, equivalently, such that $\langle 
 R_{\alpha}, F_{\beta} \rangle = \delta_{\alpha, \beta}$. 

 The \emph{immaculate basis} of $\textsf{NSym}$ may be defined by analogy with the Pieri rule in \eqref{originalPieri}, with 
\begin{equation}\label{immaculatePieri}
 \mathfrak{S}_{\alpha} H_{s} = \sum_{\beta} \mathfrak{S}_{\beta}, 
\end{equation}
 where the sum in \eqref{immaculatePieri} is over all compositions $\beta \vDash |\alpha| + s$
 that differ from $\alpha$ by an immaculate horizontal strip, i.e., so that 
 $\alpha_{j} \leq \beta_{j}$
 for all $j \in \{ 1, 2, \ldots, \ell(\alpha) \}$, and $\ell(\beta) \leq \ell(\alpha) + 1$ \cite{BergBergeronSaliolaSerranoZabrocki2014}. This basis also has a combinatorial interpretation in terms of  tableaux. An \emph{immaculate tableau} of shape $\alpha$ and content $\beta$ is a tableau $T$ of the specified shape such that the number of 
 labels in $T$ equal to $i$ is $\beta_i$ for $i \in \{ 1, 2, \ldots, \ell(\beta) \}$ and such that the first column of $T$ is strictly increasing and 
 such that the rows of $T$ are weakly increasing. Let $K_{\alpha, \beta}^{\mathfrak{S}}$ denote the number immaculate tableaux of 
 shape $\alpha$ and content $\beta$. For the case whereby $\beta = (1^{|\alpha|})$, an immaculate tableau of shape $\alpha$ and content 
 $\beta$ is said to be \emph{standard}, and we let $g^{\alpha}$ denote the number of standard immaculate tableaux of shape $\alpha$. 

\begin{example}
 In contrast to Example \ref{exampleSY}, for $\alpha = (1,2,2)$, we find that there are $g^{\alpha} = 3$ standard immaculate tableaux of 
 the given shape, as below: $$ \ytableausetup {centertableaux,boxsize=normal,notabloids}
 \begin{ytableau} 3 & 5 \\ 2 & 4 \\ 1 \end{ytableau} \ \ \ \ \ \ \ 
 \ytableausetup {centertableaux,boxsize=normal,notabloids}
 \begin{ytableau} 3 & 4 \\
 2 & 5 \\ 1 \end{ytableau}
 \ \ \ \ \ \ \ 
 \ytableausetup {centertableaux,boxsize=normal,notabloids}
 \begin{ytableau} 4 & 5 \\ 
 2 & 3 \\ 1 \end{ytableau}. $$ 
\end{example}

 For compositions $\alpha$ and $\beta$, 
 the \emph{lexicographic order} $\leq_{\ell}$ 
 may be defined recursively so that 
 $\alpha \geq_{\ell} \beta$ if
 $\alpha_1 > \beta_1$, or $\alpha_1 = \beta_1$
 and $(\alpha_2, \ldots, \alpha_{\ell(\alpha)}) \geq_{\ell} 
 (\beta_{2}, \ldots, \beta_{\ell(\beta)})$, 
 and similarly for words (over a totally ordered set). 
The iterative application of the Pieri rule in 
 \eqref{immaculatePieri} yields 
\begin{equation}\label{Htoimmaculate} 
 H_{\beta} = \sum_{\alpha \geq_{\ell} \beta} 
 K_{\alpha, \beta}^{\mathfrak{S}} \mathfrak{S}_{\alpha},
\end{equation} 
 for the number $K_{\alpha, \beta}$ of immaculate tableaux of shape $\alpha$ and content $\beta$, i.e., the number of tableaux of shape 
 $\alpha$ with a strictly increasing first column and with weakly increasing rows and with content $\beta$, i.e., so that the number of cells 
 with label $i$ is equal to $\beta_i$ for $i \in \{ 1, 2, \ldots, \ell(\beta) \}$. The expansion in \eqref{Htoimmaculate} provides a key to our 
 construction of an analogue of immaculate functions based on our lexical generalization of immaculate tableaux. 

 For general background material on Hopf algebras, we highlight standard references on Hopf algebras 
 \cite{DascalescuNastasescuRaianu2001,MilnorMoore1965}, and, for the Combinatorial Hopf Algebra
 structures on $\textsf{Sym}$, $\textsf{QSym}$, and $\textsf{NSym}$, we refer to the seminal paper by Aguiar et al.\ on CHAs \cite{AguiarBergeronSottile2006} and related references. 

\section{Lexical tableaux}\label{sectionLexical}
 Letting $ \left\{ \begin{smallmatrix} n \\ k \end{smallmatrix} \right\} $ denote the Stirling number of the second kind 
 giving the number of ways of partitioning the set $\{ 1, 2, \ldots, n \}$ into $k$ blocks, we obtain the following analogue of \eqref{krows}. 

\begin{theorem}\label{thm:snk}
 For $1 \leq k \leq n$, we have 
\begin{equation}\label{displaysnk}
 \sum_{\substack{\alpha \vDash n\\ \ell(\alpha) = k}} g^{\alpha} = 
 \left\{ \begin{matrix} 
 n \\ k 
 \end{matrix} \right\}. 
\end{equation} 
\end{theorem}

\begin{proof} 
 We construct a bijection $f$ between standard immaculate tableaux of size $n$ with $k$ rows to set partitions of $\{ 1, 2, \ldots, 
 n \}$ with $k$ blocks as follows. Given a standard immaculate tableau $T$, we let $f(T)$ be the set partition $B_1 / B_2 / \cdots / 
 B_k$ where $i \in B_j$ if $i$ is in Row $j$ of $T$. Note that the order of the blocks is irrelevant since we are working with unordered set 
 partitions. The inverse map $f^{-1}$ takes a set partition $\pi = B_1 / B_2 / \cdots / B_k$ and constructs a standard immaculate 
 tableaux $f^{-1}(\pi)$ of shape $\text{sort}(|B_{j_1}|, |B_{j_2}|, \ldots, |B_{j_k}|)$ where $\min(B_{j_1}) < \min(B_{j_2}) < \cdots < 
 \min(B_{j_k})$, and where the entries of row $i$ are exactly the entries of $B_{j_i}$ sorted into increasing order. 
\end{proof}

 Theorem \ref{thm:snk} illustrates how immaculate tableaux are useful and natural combinatorial objects, and the
 right-hand evaluation in \eqref{displaysnk} leads us to consider what would be appropriate as an analogue of 
 \eqref{displaysnk} based on immaculate-like tableaux and variants or generalizations of $ \left\{ \begin{smallmatrix} 
 n \\ k \end{smallmatrix} \right\}$. In this direction, the unsigned Stirling number $|s(n,k)|= {n \brack k}$ of the first kind is equal to the number of permutations of $n$ elements with $k$ disjoint cycles. Given a permutation $\sigma$ of $\{1,2,\ldots, n\}$ with $k$ disjoint cycles, we denote its cycle decomposition $c_{\sigma, 1}, c_{\sigma, 2}, \ldots, c_{\sigma, k}$, with each cycle written as a tuple of elements in $\{1,2,\ldots, n\}$. We order these tuples from least to greatest based on their minimal element.

 To construct a new family of tableaux based on an analogue of Theorem \ref{thm:snk} with 
 ${n \brack k} $ in place of 
 $ \left\{ \begin{smallmatrix} 
 n \\ k 
 \end{smallmatrix} \right\}$, 
 we employ the concept of a \emph{cyclic shift}, as seen in the work of Adin et al.\ \cite{AdinGesselReinerRoichman2021}. Given a word $w = w_1 w_2 \cdots w_{\ell(\ell)}$, 
 a cyclic shift of $w$ is given by $w^{(i)} = w_{i+1} w_{i+2} \cdots
 w_{\ell(w)} w_1 w_2 \cdots w_{i}$ for any $i \in [ \ell(w) 
 ] = \{ 1, 2, \ldots, \ell(w) 
 \}$. Let $[\vec{w}]$ denote the set of cyclic shifts of $w$. Define a \emph{necklace word} to be a word that is lexicographically minimal among all of its cyclic shifts.

\begin{definition}
 We define a \emph{lexical tableaux} of shape $\alpha \vDash n$ and type (or content) $\beta$ as a filling of the diagram
 of $\alpha$ such that $i$ appears exactly $\beta_i$ times, the entries in the first column are strictly increasing, and the word
 $w_i$ formed by the entries of row $i$ (in order) is a necklace word.
\end{definition}

 We refer to lexical tableaux of type $\beta = (1)^n$ as \emph{standard}. 

\begin{theorem}\label{theoremltstirling}
 Let $lt_{\alpha}$ be the number of  standard  lexical tableaux of shape $\alpha$. Then 
\begin{equation}\label{displayltstirling}
 \sum_{\substack{ \alpha \vDash n,\\ \ell(\alpha)=k}} lt_{\alpha} = {n \brack k}. 
\end{equation} 
\end{theorem}

\begin{proof} 
 Let $T$ be a standard lexical tableau of size $n$ with $k$ rows. We map this to a permutation $\sigma$ of $[n]$ by setting $c_{\sigma, i}$ to be the tuple formed by the entries in row $i$ of $T$. This map is a bijection, 
 where the inverse map takes a permutation with cycle decomposition $c_{\sigma,1}c_{\sigma, 2} \cdots c_{\sigma, k}$ to a standard lexical tableau where the entries in row $i$ are given by the entries in $c_{\sigma, i}$, in the same order.
\end{proof}

\begin{example}
 For the ${4\brack2}=11$ case,   the standard lexical tableaux with $4$ blocks and $2$ rows are as below. \vspace{2mm}
\[
 \begin{ytableau}
 2 & 3 & 4\\
 1
 \end{ytableau} \qquad
 \begin{ytableau}
 2 & 4 & 3\\
 1
 \end{ytableau} \qquad
 \begin{ytableau}
 2\\
 1 & 3 & 4
 \end{ytableau} \qquad
 \begin{ytableau}
 2\\
 1 & 4 & 3
 \end{ytableau} \qquad
 \begin{ytableau}
 3\\
 1 & 2 & 4
 \end{ytableau} \qquad
 \begin{ytableau}
 3\\
 1 & 4 & 2
 \end{ytableau} \] \vspace{-4mm}
 
 \[
 \begin{ytableau}
 4\\
 1 & 2 & 3
 \end{ytableau} \qquad
 \begin{ytableau}
 4\\
 1 & 3 & 2
 \end{ytableau}\qquad
 \begin{ytableau}
 3 & 4\\
 1 & 2 
 \end{ytableau} \qquad
 \begin{ytableau}
 2 & 4\\
 1 & 3 
 \end{ytableau} \qquad
 \begin{ytableau}
 2 & 3\\
 1 & 4
 \end{ytableau} 
 \] \vspace{2mm}
 
 This illustrates how lexical tableaux generalize standard immaculate tableaux, since for the $ \left\{ \begin{smallmatrix} 4 \\ 2 \end{smallmatrix} \right\} = 7$ case, 
 the standard immaculate tableaux with $4$ blocks and $2$ rows are as below. \vspace{2mm}
\[
 \begin{ytableau}
 2 & 3 & 4\\
 1
 \end{ytableau} \qquad
 \begin{ytableau}
 2\\
 1 & 3 & 4
 \end{ytableau} \qquad
 \begin{ytableau}
 3\\
 1 & 2 & 4
 \end{ytableau} \qquad
 \begin{ytableau}
 4\\
 1 & 2 & 3
 \end{ytableau} \qquad
 \begin{ytableau}
 3 & 4\\
 1 & 2 
 \end{ytableau} \qquad
 \begin{ytableau}
 2 & 4\\
 1 & 3 
 \end{ytableau} \qquad
 \begin{ytableau}
 2 & 3\\
 1 & 4
 \end{ytableau} 
 \] \vspace{2mm}
\end{example}

 Two tableaux $T_1$ and $T_2$ of a given shape with $\ell$ rows are said to be \emph{row-equivalent}, if the $i^{\text{th}}$ rows of 
 $T_1$ and $T_2$ contain the same labels for all $i \in \{ 1, 2, \ldots, \ell \}$. This gives rise to an equivalence relation $\sim$, writing 
 $T_1 \sim T_2$ if $T_1$ and $T_2$ are row-equivalent. A \emph{tabloid} may then be defined as the equivalence class associated with 
 $\sim$ of a standard Young tableau, and may be denoted with any member of the corresponding equivalence class, with the notational 
 convention whereby vertical bars are removed. 

\begin{example}
 The expressions $$ \ytableausetup {boxsize=normal,tabloids} \ytableaushort{ 3,12 } \hfill \ \ \ \text{and} \ \ \ 
 \ytableausetup {boxsize=normal,tabloids} 
 \ytableaushort{ 3,21 } \hfill $$
 denote the same tabloid of shape $(2, 1)$. 
\end{example}

 Tabloids, as defined above, play roles of basic importance in the representation theory of the symmetric group, with reference to Sagan's 
 classic text \cite[\S2.1]{Sagan2001}. The above definition of a lexical tableau may be seen as providing an analogue for immaculate 
 tableaux of tabloids, and this is formalized below. 

\begin{theorem}\label{theoremclasses}
 The number of equivalence classes, with respect to the row-equivalence relation $\sim$, on standard lexical tableaux of shape $\alpha$ 
 is equal to the number of standard immaculate tableaux of shape $\alpha$.
\end{theorem}

\begin{proof}
 By taking a standard immaculate tableau $T$ of a given shape $\alpha$, the standard lexical tableaux of the same shape are obtained 
 by permuting labels within each row and to the right of the first column. This forms a bijection giving the desired result, by taking $T$ as 
 the representative of the equivalent class $[T]_{\sim}$.
\end{proof}

 Theorem \ref{theoremclasses}, together with how \eqref{displayltstirling} provides a natural compantion to the identity in 
 \eqref{displaysnk} involving standard immaculate tableaux, motivate  the problem of constructing an analogue of immaculate and dual 
 immaculate functions with the use of lexical tableaux in place of immaculate tableaux. This forms the main purpose of our paper and is 
 motivated by the importance of immaculate functions within many different areas of algebraic combinatorics. 
 
 In our construction of bases of $\textsf{QSym}$ and $\textsf{NSym}$ via lexical tableaux, we require properties on the enumeration of 
 lexical tableaux, and hence the material in Section \ref{subsectionEnumeration} below. 

\subsection{On the enumeration of lexical tableaux}\label{subsectionEnumeration}
 Consider a mulitset $\mathcal{B} =\{a_1^{n_1}, a_2^{n_2}, \ldots, a_k^{n_k} \}$. The number of necklace words with characters corresponding exactly to the elements in $\mathcal{B}$ is given by \[N(\mathcal{B}) = \frac{1}{|\mathcal{B}|} \sum_{d|\gcd(n_1, \ldots, n_k)} \binom{|\mathcal{B}|/d}{n_1/d, \ldots, n_k/d} \varphi(d),\]
where $\varphi$ is Euler's totient function \cite{gilbert1961symmetry}. Let $IT_{\alpha,\beta}$ be the set of immaculate tableaux of shape $\alpha$ and type $\beta$. Let $\mathcal{R}^T_i$ be the multiset of entries in row $i$ of an immaculate tableau $T$. 
 Then, we have \[ K^{\lx}_{\alpha, \beta} = \sum_{T \in IT_{\alpha, \beta}} N(\mathcal{R}^T_1)N(\mathcal{R}^T_2) \cdots N(\mathcal{R}^T_{\ell(\alpha)}).\]

 We can also count standard lexical tableaux with methods coming from the study of immaculate tableaux. Given a cell $c = (i,j)$ in a composition 
 diagram of $\alpha$, a hook of $c$ in $\alpha$, denoted $h_{\alpha}(c),$ is defined to be the number of cells below and to the right of $c$ if $c$ is in the first column, and the number of cells weakly to the right of $c$ in the same row otherwise. That is, if $j=1$, we have 
 $h_{\alpha}(c) = \alpha_{i} + \alpha_{i+1} + \cdots \alpha_{k}$. If $j>1$ then $h_{\alpha}(c)=\alpha_i-j+1$. Berg et al.\ 
 \cite{BergBergeronSaliolaSerranoZabrocki2014}, 
 proved that the number of standard immaculate tableaux of shape $\alpha$, denoted here as $K^{\mathfrak{S}}_{\alpha, 1^n}$, is 
 equal to \[ K^{\mathfrak{S}}_{\alpha, 1^n} = \frac{n!}{\prod_{c \in \alpha}h_{\alpha}(c)}.\]
 This leads us to the following formula for the number of standard lexical tableaux.

\begin{theorem} 
 Let $\alpha \vDash n$. The number of standard lexical tableaux of shape $\alpha$ is given by \[ K^{\lx}_{\alpha, 1^n} = \frac{n! \prod_{i \in [\ell(\alpha)]} (\alpha_i-1)!}{\prod_{c \in \alpha} h_{\alpha}(c)}.\]
\end{theorem}

\begin{proof}
For each standard immaculate tableau of shape $\alpha$, we can generate $\prod_{i \in [k]}(\alpha_i - 1)!$ unique lexical tableaux by permuting the entries within each row, excluding those in the first column. No two lexical tableau generated this way can be the same, and every lexical tableau is associated with some immaculate tableau in this way, so all will be generated.
\end{proof}

 The following result is key in relation to the 
 unitriangularity of transition matrices we later require. 

\begin{theorem}\label{theorem:lexical_coeff}
 Given a composition $\alpha \vDash n$, we have $K^{\lx}_{\alpha, \alpha}=1$ and, if $\alpha \leq_{\ell} \beta$, then $K^{\lx}_{\alpha, \beta}=0$.
\end{theorem}

\begin{proof}
 Consider an empty diagram of shape $\alpha = (\alpha_1, \ldots, \alpha_k)$ that we want to fill as a lexical tableau of type $\alpha$. If 
 there is a $1$ anywhere in a given row, then the first entry of that row must be a $1$. Since the first column is strictly increasing, there
 must be a $1$ in the first row, and there must not be any $1$s in the following rows. So all $\alpha_1$ of the $1$s in the tableaux must 
 go in the top row, and they fill it completely. Next, we need to fill $\alpha_2$ cells with $2$. Here, $2$ will be the lowest entry in any of 
 the rows (other than the first), so any row with a $2$ must have a $2$ as its first entry. Since the first column is strictly increasing, 
 the second row must have a $2$ as its first entry, and no other row may contain any $2$'s. Therefore, we fill the second row entirely 
 with $2$'s, which uses all $\alpha_2$ that we seek to use. Continuing this pattern, we see that the only way to construct a lexical tableau 
 of shape $\alpha$ and type $\alpha$ is to fill each of the $\alpha_i$ cells of row $i$ with $i$'s, and thus  $K_{\alpha,\alpha} = 1$. 
 
 Next, consider some $\alpha \leq_{\ell} \beta$, meaning there exists some $j$ such that $\alpha_1 = \beta_1, \alpha_2 = \beta_2, \ldots, 
 \alpha_{j-1}=\beta_{j-1},$ and $\alpha_j<\beta_j$. By way of contradiction, suppose that there exists a lexical tableau $T$ of shape $\alpha$ and type $\beta$. Using a similar argument, relative to our preceding argument, the first $j-1$ rows of $T$ are necessarily filled entirely with the integer that matches their row index. That is, if $i < j$ then row 
 $i$ is has exactly $\alpha_i = \beta_i$ each filled with an $i$. Next, we will fill in the $\beta_j$ instances of $j$. Given the current filling, 
 any row with a $j$ must have a $j$ as its first entry. Since the first column is strictly increasing, row $j$ must contain all of these $j$'s. 
 However, there are only $\alpha_j$ cells in row $j$ and we need to place $\beta_j > \alpha_j$ instances of $j$. Thus, we cannot create a 
 lexical tableau of shape $\alpha$ and type $\beta$, meaning $K^{\lx}_{\alpha, \beta}=0$ if $\alpha \leq_{\ell} \beta$. 
\end{proof}

\begin{table}[h!]\label{tab:trans}
 \centering
 \begin{tabular}{|c|c|c|c|c|c|c|c|c|} \hline
 $K^{\lx}_{\alpha. \beta}$ & (4) & (3,1) & (2,2) & (2,1,1) & (1,3) & (1,2,1) & (1,1,2) & (1,1,1,1) \\ \hline
 (4) & 1 & 1 & 2 & 3 & 1 & 3 & 3 & 6 \\ \hline
 (3,1) & 0 & 1 & 1 & 2 & 1 & 3 & 3 & 6 \\ \hline
 (2,2) & 0 & 0 & 1 & 1 & 1 & 2 & 2 & 3 \\ \hline
 (2,1,1) & 0 & 0 & 0 & 1 & 0 & 1 & 1 & 3 \\ \hline
 (1,3) & 0 & 0 & 0 & 0 & 1 & 1 & 1 & 2 \\ \hline
 (1,2,1) & 0 & 0 & 0 & 0 & 0 & 1 & 1 & 2 \\ \hline
 (1,1,2) & 0 & 0 & 0 & 0 & 0 & 0 & 1 & 1 \\ \hline
 (1,1,1,1) & 0 & 0 & 0 & 0 & 0 & 0 & 0 & 1 \\ \hline
 
 \end{tabular}
 \caption{Values of $K^{\lx}_{\alpha, \beta}$ for $\alpha, \beta \vDash 4$.}
 \label{tab:placeholder}
\end{table}

For our next results, we extend the standardization map on immaculate tableaux to lexical tableaux. If $T$ is a lexical tableau of shape $\alpha$, then $std(T)=S$ will be the standard lexical tableau of shape $\alpha$ created as follows. Read through the entries in $T$ starting first with those equal to 1, then 2, etc. Among all entries of the same value in $T$, read from left to right and top to bottom. Replace entries in the order they are read, starting with 1 and increasing each entry. Note that, like with immaculate tableaux, no two lexical tableaux of the same shape and type can have the same standardization. 

 Given a standard lexical tableau $S$ of shape $\alpha$, let $\Max(S)$ be the set of maximal elements in terms of  the refinement ordering   on  the set of compositions $\gamma$ for which a lexical tableau of shape $\alpha$ and type $\gamma$ exists and standardizes to $S$. Let $J^{\lx}_{\alpha,\beta}$ be the number of standard lexical tableaux $S$ of shape $\alpha$ such that $\beta \in \Max(S)$.

\ytableausetup {boxsize=normal,notabloids}
\begin{example}
 The standard lexical tableau  $S$ has $Max_{\succeq}(S)=\{T_1, T_2\}$, where \[ S = \begin{ytableau}
 1 & 2 & 4 & 3
 \end{ytableau}\ , \qquad \quad T_1 = \begin{ytableau}
 1 & 2 & 3 & 2
 \end{ytableau}\ , \quad T_2 = \begin{ytableau}
 1 & 1 & 3& 2
 \end{ytableau}. \] 
\end{example}

\begin{lemma}\label{lem:ref_lex}
 Let $\beta \succeq \gamma$ and $\alpha$ be compositions of $n$. If there exists a lexical tableau $T$ of shape $\alpha$ and type $\beta$, then there exists a lexical tableau $R$ of shape $\alpha$ and type $\gamma$ such that $std(T)=std(R)$.
\end{lemma}
\begin{proof}
 Let $T$ be a lexical tableau of shape $\alpha$ and type $\beta$, and let $\gamma \preceq \beta$. Let $R$ be the tableau given by filling a diagram of shape $\alpha$ with the entries with $\gamma$ in the order the slots are numbered in $std(T)$. We will show that $R$ is a lexical tableau.

 Assume for contradiction that $R$ is not a lexical tableau, meaning that there is some row $i$ that is not filled by a necklace word. 
 So, row $i$ of $R$ is  filled with a word of the form $w \cdot v$ where $v \cdot w$ is a necklace word and $v \cdot w \leq_{\ell} w \cdot v$. Let $|w|=j$ and $|v| = \alpha_i-j$. Then, row $i$ of $T$ is filled by some word of the form $w' \cdot v'$ where $|w'|=j$ and $|v'|=\alpha_i - j$. Finally, let $u = u_1u_2 \cdots u_{\alpha_i}$ be row $i$ of $std(T)$ and thus $std(R)$ as well. 

 Since $v \cdot w \leq_{\ell} w \cdot v$, it must be that $v \leq_{\ell} w$. We have two possible cases. In case (1), we have $v_1 = w_1$ through $v_{r-1} = w_{r-1}$ and $v_{r} < w_{r}$ for some $r \in [min(j, \alpha_i -j)]$. As a result, it must be that $u_{j+r} < u_{r}$ by our standardization. Thus, we must have $v'_r < w'_r$ as well.
 So, we have that $w'_1 = v'_1, \ldots, w'_{t-1} = v'_{t-1}$ with $w'_{t} \not= v'_{t}$ in $T$ with $t \leq r$. However, we must also have $r \leq t$ because if $w_i = v_i$ in $R$ then we must have $w'_i = v'_i$ in $T$. Therefore, it must be that $t=r$ and so $v' \leq_{\ell} w'$. In case (2), $w_1 = v_1, \ldots, w_{\alpha_i - j} = v_{\alpha_i-j}$ where $j > \alpha_i - j$, and so $w'_1 = v'_1, \ldots, w'_{\alpha_i - j} = v'_{\alpha_i-j}$. Thus, $v' \leq_{\ell} w'$. 

 In both cases, then, we have $v' \cdot w' \preceq w' \cdot v'$, meaning row $i$ of $T$ is not a necklace word. This is a contradiction as $T$ is a lexical tableau. Thus, there must not exist any row in $R$ that is not filled by a necklace word, and so $R$ is a lexical tableau.
\end{proof}

\begin{theorem}\label{thm:K-to-J} For $\alpha, \gamma \models n \in \N$, we have 
\[ K^{\lx}_{\alpha, \gamma} = \sum_{\beta \succeq \gamma} J^{\lx}_{\alpha,\beta}.\]
\end{theorem}

\begin{proof}
Let $LT_{\alpha,\gamma}$ be the set of lexical tableaux of shape $\alpha$ and type $\gamma$. Let $SLT_{\alpha}(\succeq\gamma)$ be the set of standard lexical tableaux $S$ of shape $\alpha$ where there is some $\beta \succeq \gamma$ such that $\beta \in Max(S)$. We will show that standardization is a map between these two sets. Since $|LT_{\alpha, \gamma}|=K^{\lx}_{\alpha, \gamma}$ and $|SLT_{\alpha}(\succeq \gamma)| = \sum_{\beta \succeq \gamma} J^{\lx}_{\alpha, \beta}$, this proves the claim.

Let $T \in LT_{\alpha, \gamma}$. By definition, there must exist some $\beta \succeq \gamma$ such that $\beta \in Max(std(T)).$ Thus, $std(T) \in SLT_{\alpha}(\succeq \gamma)$. Moreover, there will be no other lexical tableaux of the same shape and type that standardize to $std(T)$,  so that  $std: LT_{\alpha, \gamma} \rightarrow SLT_{\alpha}(\succeq \gamma)$ is injective. Next, consider some $S \in SLT_{\alpha}(\succeq \gamma)$. By Lemma \ref{lem:ref_lex}, there    necessarily
  exists some lexical tableau  $R$ of shape $\alpha$ and type $\beta$ with $\beta \succeq \gamma$ and $std(R)=S$ since $S \in SLT_{\alpha}(\succeq \gamma)$. Thus $std: LT_{\alpha, \gamma} \rightarrow SLT_{\alpha}(\succeq \gamma)$ is surjective and so it is a bijection. 
\end{proof}

This identity will be crucial in our expansion of the ribbon basis into the lexical functions in the following section. 

\section{Lexical functions in $\textsf{NSym}$ and $\textsf{QSym}$}\label{sectionqsym}
 Let $LT_{\alpha}$ denote the set of all  lexical tableaux of shape $\alpha$. Given $T \in LT_{\alpha}$, let $type(T)$ denote the
 type of $T$. If $type(T)=\beta= (\beta_1, \beta_2, \ldots, \beta_k)$, we associate $T$ with the monomial $x^T = x_1^{\beta_1} 
 x_2^{\beta_2} \cdots x_k^{\beta_k}$. Additionally, let $K^{\lx}_{\alpha,\beta}$ denote the number of lexical tableaux of shape 
 $\alpha$ and type $\beta$.

\begin{definition}
 For a composition $\alpha \vDash n$, define the dual lexical function by \[ \lx^*_{\alpha} = \sum_{T \in LT_{\alpha}} x^T.\]
\end{definition}

Note that we define these as the \emph{dual} lexical functions so that we may call their duals simply the lexical functions. This is consistent with the immaculate and dual immaculate functions of $\textsf{NSym}$ and $\textsf{QSym}$.

 Given a lexical tableau $T$ with entries given by the set $N$ with $|N|=n$, let $pack(T)$ be the lexical tableau that is
 obtained by replacing the entries from $N$ in $T$ with entries in $[n]$ according to the unique order-preserving bijetion between $N$ and $[n]$. We call a lexical tableau \emph{packed} if $pack(T) = T$, or in other words, its entries correspond exactly to the elements of the set $[n]$. These are exactly the lexical tableaux whose type is given by a strong composition, as opposed to a weak composition.

\begin{theorem}\label{theoremlxstarM}
The dual lexical functions have a positive, uni-triangular expansion in terms of the monomial quasisymmetric functions as \[ \lx_{\alpha}^* = \sum_{\beta \vDash n} K^{\lx}_{\alpha, \beta} M_{\beta}.\] 
\end{theorem}

\begin{proof} 
Observe that, for a packed lexical tableau $T_p$ of type $\beta$, we have $\sum_{pack(T) = T_p} x^T = M_{\beta}$. Then, \[\sum_{T \in LT_{\alpha}} x^T = \sum_{\substack{T \in LT_{\alpha},\\ pack(T) = T}} M_{type(T)} = \sum_{\beta \vDash 
 n} K^{\lx}_{\alpha,\beta} M_{\beta}.\] By Theorem \ref{theorem:lexical_coeff}, the transition matrix from the $M$ to $\lx$ will be unitriangular when the indices are arranged in lexicographic order.
\end{proof}

\begin{example}

We have the expansion \[\lx^*_{(3,1)} = M_{(3,1)} + M_{(2,2)} + 2M_{(2,1,1)} + M_{(1,3)} + 3M_{(1,2,1)} + 3M_{(1,1,2)} + 6M_{(1,1,1,1)},\]

corresponding to the lexical tableaux

\[ \begin{ytableau}
 2\\
 1 & 1 & 1 
\end{ytableau} \quad \begin{ytableau}
 2\\
 1 & 1 & 2 
\end{ytableau} \quad \begin{ytableau}
 3\\
 1 & 1 & 2 
\end{ytableau} \quad \begin{ytableau}
 2\\
 1 & 1 & 3
\end{ytableau} \quad
\begin{ytableau}
 2\\
 1 & 2 & 2 
\end{ytableau} 
\] \vspace{-2mm}

\[
\begin{ytableau}
 2\\
 1 & 3 & 3
\end{ytableau} \quad
\begin{ytableau}
 3\\
 1 & 2 & 3
\end{ytableau} \quad
\begin{ytableau}
 3\\
 1 & 3 & 2
\end{ytableau}
\quad \begin{ytableau}
 3\\
 1 & 2 & 2
\end{ytableau} \quad \begin{ytableau}
 2\\
 1 & 2 & 3
\end{ytableau} \quad \begin{ytableau}
 2\\
 1 & 3 & 2
\end{ytableau}
\] \vspace{-2mm}

\[\begin{ytableau}
 4\\
 1 & 2 & 3
\end{ytableau} \quad \begin{ytableau}
 4\\
 1 &3 & 2
\end{ytableau} \quad \begin{ytableau}
 3\\
 1 & 2 & 4
\end{ytableau} \quad
\begin{ytableau}
 3\\
 1 & 4 & 2
\end{ytableau} \quad
\begin{ytableau}
 2\\
 1 & 3 & 4
\end{ytableau} \quad
\begin{ytableau}
 2\\
 1 & 4 & 3
\end{ytableau}.
\] \vspace{0mm}
 
\end{example}

From Theorem \ref{theoremlxstarM} we have the following result.

\begin{corollary}
 The set $\{\lx^*_{\alpha} : \alpha \vDash 
 n\}$ is a basis for $\emph{\textsf{QSym}}_n$ and $\bigcup_{n \geq 0}\{ \lx^*_{\alpha} : \alpha \vDash
 n\}$ is a basis for $\emph{\textsf{QSym}}$.
\end{corollary}

\begin{remark}
 Observe that if $\alpha \vDash n$ and $\alpha_i \leq 2$ for all $i \in [\ell(\alpha)]$, then $\lx^*_{\alpha} = \mathfrak{S}^*_{\alpha}$. 
 This may be seen using the property such that a necklace word of length two is always a weakly increasing word. Thus, the lexical 
 tableaux of shape $\alpha$ where $\alpha_i \leq 2$ for all $i \in \ell(\alpha)$ are exactly the immaculate tableaux of shape $\alpha$.
\end{remark}

Now we define the dual basis of the lexical functions in $\textsf{NSym}$.

\begin{definition}\label{definitionlexinner}
 Define the lexical basis of $\textsf{NSym}$ to be the unique basis $\bigcup_{n \geq 0}\{\lx_{\alpha} : \alpha \vDash 
 n\}$ such that 
 \begin{equation}\label{innerforlex}
 \langle \lx_{\alpha}, \lx^{*}_{\beta} \rangle = \delta_{\alpha,\beta},
 \end{equation} 
 for all compositions $\alpha,\beta \vDash 
 n$ for all $n \geq 0$.
\end{definition}
 From the inner product relation in \eqref{innerforlex}, we obtain the below analogue of the $H$-to-$\mathfrak{S}$ expansion 
 formula in \eqref{Htoimmaculate} due to Berg et al.\ \cite{BergBergeronSaliolaSerranoZabrocki2014}. 

\begin{corollary}\label{Htolex}
 For $\beta \vDash n$, we have the expansion \[H_{\beta} = \sum_{\alpha 
 \geq_{\ell} \beta} K^{\lx}_{\alpha, \beta} \lx_{\alpha}.\]
\end{corollary}

\begin{proof}
 This follows from Theorem \ref{theoremlxstarM} together with the duality relation in Definition \ref{definitionlexinner}. 
\end{proof}

  Using  Theorem \ref{thm:K-to-J}, we can also give a positive expansion of the ribbon basis of $\textsf{NSym}$ into the lexical functions.

\begin{theorem}\label{Rtolex}
 The ribbon noncommutative symmetric functions expand into lexical functions as 
 \[R_{\beta} = \sum_{\alpha \models n} J^{\lx}_{\alpha, \beta} \lx_{\alpha}.\]
\end{theorem}

\begin{proof}
 By Theorem \ref{theoremlxstarM}, we have $H_{\gamma} = \sum_{\alpha \models n} K^{\lx}_{\alpha, \gamma} \lx_{\alpha}$. Thus, using Theorem \ref{thm:K-to-J},
 \[H_{\gamma} = \sum_{\alpha \models n} K^{\lx}_{\alpha,\gamma} \lx_{\gamma} = 
 \sum_{\alpha \models n} \sum_{\beta \succeq \gamma} J^{\lx}_{\alpha, \beta} \lx_{\beta} = \sum_{\beta \succeq \gamma} \left( \sum_{\alpha \models n} J^{\lx}_{\alpha, \beta}\lx_{\alpha} \right)\]
 Given that $H_{\gamma} = \sum_{\beta \succeq \gamma} R_{\beta}$, it must be that $R_{\beta} = \sum_{\alpha \models n} J^{\lx}_{\alpha, \beta}\lx_{\alpha}$.
\end{proof}

 Theorem \ref{Rtolex}  also yields the dual expansion in $\textsf{QSym}$.

\begin{corollary}
 The dual lexical functions have a positive expansion in terms of the fundamental quasysymmetric functions as \[ \lx^*_{\alpha}= \sum_{\beta \models n} J^{\lx}_{\alpha,\beta}F_{\beta}.\]
\end{corollary}

\begin{example}
 We have the expansion \[ \lx^*_{(4)} = F_{(1,1,2)} + 2F_{(1,2,1)} + F_{(2,1,1)} + F_{(2,2)} + F_{(4)}.\]
\end{example}

 The difficulties, from both computational and combinatorial points of view, 
 associated with problems related to the evaluation of $\lx$-basis elements, are reflected by how
 the $\lx$-basis
 does not satisfy any Pieri rule with $\{ 0, 1 \}$-coefficients, as illustrated below. 

\begin{example}
  Consider the expansion 
\begin{multline*}
 \lx_{{{\tikztableausmall{ {X}, {X, X}} }}} 
 H_{{{\tikztableausmall{ {\boldentry X, \boldentry X}} }}} = 
 \lx_{{{\tikztableausmall{ {\boldentry X, \boldentry X}, {X}, {X, X}} }}} + \lx_{{{\tikztableausmall{ {\boldentry X}, {X, \boldentry X}, {X, X}} }}} + \lx_{{{\tikztableausmall{ {X, \boldentry X, \boldentry X}, {X, X}} }}} + \lx_{{{\tikztableausmall{ {\boldentry X}, {X}, {X, X, \boldentry X}} }}} + \\ 
 \lx_{{{\tikztableausmall{ {X, \boldentry X}, {X, X, \boldentry X}} }}} + 
 4 \lx_{{{\tikztableausmall{ {X}, {X, X, \boldentry X, \boldentry X}} }}} + 
 4 \lx_{{{\tikztableausmall{ {X, X, \boldentry X, \boldentry X, \boldentry X}} }}}.  
\end{multline*} 
\end{example}

 Informally, since lexical tableaux generalize immaculate tableaux, the lexical basis may be seen as more complicated  or intractable than  
 the immaculate   basis, and the fact that the lexical basis does not satisfy any Pieri rule may be seen as representative of this. For example,  
 one might  hope to obtain an $E$-to-$\lx$ expansion formula using Corollary \ref{Htolex}, but the known   $E$-to-$\mathfrak{S}$ 
 expansion formula relies on a Pieri rule for products of the form  $\mathfrak{S}_{\alpha} E_{s}$,  but, as above,    products of the form  
 $\lx_{\alpha} E_{s}$ do not satisfy a Pieri rule.   In a similar spirit, the intractable nature of the $\lx$-basis is such that it does not seem to be  
  feasible to construct   this basis using a Jacobi--Trudi-like or determinantal      rule or with analogues of Bernstein operators, in contrast to   
   the immaculate basis. 

 Given an open problem concerning the immaculate basis, one might consider a corresponding problem for the lexical basis, in the hope that the use of lexical basis elements could shed light on the original problem, and this provides a main source of 
 motivation concerning the study of the $\lx$-basis. In this direction, the problem of determining cancellation-free formulas for the 
 antipode $S = S_{\textsf{NSym}}$ mapping of $\textsf{NSym}$ evaluated at immaculate basis elements remains open, despite past progress 
 on this problem \cite{BenedettiSagan2017,Campbell2023}. Since progress on this problem has been made for cases given by specific 
 families of composition shapes, we consider the problem 
 of determining cancellation-free formulas for lexical functions indexed by the same composition shapes. 

\subsection{Antipodes of lexical basis elements}
 In their seminal paper on antipodes and involutions, Benedetti and Sagan \cite{BenedettiSagan2017} used a bijective approach toward 
 obtaining cancellation-free formulas for expanding $S(\mathfrak{S}_{\alpha})$ in the $\mathfrak{S}$-basis, for the cases whereby 
 $\alpha$ is a hook or consists of two rows. Subsequently, cancellation-free formulas for expanding $S(\mathfrak{S}_{\alpha})$ in the 
 $R$-basis were determined for the cases whereby $\alpha$ is a rectangle or certain products of rectangles \cite{Campbell2023}, and this 
 was later generalized by Allen and Mason \cite{AllenMason2025}. The $S(\mathfrak{S}_{\alpha})$-to-$R$ expansion formulas relied on 
 the Jacobi--Trudi formula for the immaculate basis, and the known formula for the antipode of an immaculate-hook also relied on the 
 Jacobi--Trudi formula for the $\mathfrak{S}$-basis, but the $\lx$-basis does not seem to satisfy any such determinantal formula. This 
 leads us to consider the problem of evaluating $S(\lx_{\alpha})$ for the case whereby $\alpha$ is a two-rowed composition. 

\begin{example}
 From the expansion 
\begin{equation}\label{lx24}
 \lx_{24} = H_{24} - H_{33} - H_{42} + H_{51} + 3 H_{6}, 
\end{equation}
 we obtain a cancellation-free formula for $S(\lx_{24})$ by applying $S$ to both sides of \eqref{lx24} and using the property that $S$ is an 
 anti-homomorphism together with the relation $S(H_{n}) = (-1)^n E_n$. 
\end{example}

 With regard to the following lemma, we adopt the notational convention whereby: For a composition $\alpha$, the concatenation 
 $\alpha \cdot (0)$ may be identified with $\alpha$. 

\begin{theorem}\label{mainccoeff}
 For positive integers $a$ and $b$, define $$c^{(a, b)}_{1} := -K^{\lx}_{(a+1,b-1), (a, b)}$$ and then recursively define $$ c^{(a, 
 b)}_{i} := -\left( K^{\lx}_{(a+i,b-i), (a, b)} + \sum_{j=1}^{i-1} c^{(a, b)}_{j} K^{\lx}_{(a+i,b-i),(a+j,b-j)} \right) $$ for all 
 possible indices $i$. Then 
\begin{equation}\label{lxtworowexp} 
 \lx_{(a, b)} = H_{(a, b)} 
 + c_{1}^{(a, b)} H_{(a+1, b-1)} 
 + c_{2}^{(a, b)} H_{(a+2,b-2)}
 + \cdots + c_{b}^{(a, b)} H_{(a+b)}. 
\end{equation}
\end{theorem}

\begin{proof}
 From the condition whereby lexical tableaux are required to be strictly increasing in the first column, we may deduce, from Corollary 
 \ref{Htolex}, that each $\lx$-term in the expansion of $H_{(a, b)}$ is indexed by a composition of length not exceeding $2$. Moreover, 
 from the triangularity of the transition matrices between the $\lx$- and $H$-bases given by Theorem \ref{theorem:lexical_coeff}, we have 
 that   the compositions indexing the $H$-elements in the $H$-expansion of $\lx_{(a, b)}$ are of the form $\lx_{(a+i,b-i)}$ for nonnegative 
 $i$.   The desired result then follows inductively, by expanding $H_{(a, b)}$ into the $\lx$-basis and performing an equivalent form of 
 Gaussian elimination. 
\end{proof}

 As a consequence, we obtain the antipode formula 
 \begin{equation}\label{antipodeE}
 S\big(\lx_{(a, b)}\big) = (-1)^{a+b} \big( E_{(b, a)} 
 + c_{1}^{(a, b)} E_{(b-1, a + 1)} 
 + c_{2}^{(a, b)} E_{(b-2, a + 2)}
 + \cdots + c_{b}^{(a, b)}
 E_{(a+b)} \big), 
 \end{equation} 
 i.e., by applying the antipode map $S$ to both sides of the expansion in \eqref{lxtworowexp}. In view of Benedetti and Sagan's bijective 
 approach toward the cancellation-free evaluation of antipodes of the form $S(\mathfrak{S}_{(a, b)})$, this motivates the problem of 
 finding combinatorial interpretations for coefficients of the form $c_{i}^{(a, b)}$. 

 Adopting notation from the OEIS entry {\tt A047996}, let $T(n, k)$ denote the number of necklaces with $k$ black beads and $n-k$ white 
 beads. Equivalently, 
 we may define $T(n, k)$ so that 
\begin{equation}\label{Texplicit} 
 T(n, k) := \frac{1}{n}
 \sum_{d \mid n, k} \varphi(d) \binom{n/d}{k/d}. 
\end{equation} 

\begin{lemma}\label{relationT}
 The relation $$ K^{\lx}_{(a+i,b-i), (a+j,b-j)} = T(a+i,i-j) $$ holds for $j \in \{ 0, 1, \ldots, i-1 \}$. 
\end{lemma}

\begin{proof}
 Either side of \eqref{Texplicit} is equal to the number of length-$n$ binary words $w$ with $k$ entries equal to $1$ such that $w$ is 
 lexicographically less than or equal to all of the cyclic permutations of $w$. A lexical tableau of shape $(a+i,b-i)$ and content $(a+j,b-j)$ 
 is uniquely determined by its initial row, which, when read as a word, is lexicographically less than or equal to all of the cyclic 
 permutations of the same word. 
\end{proof} 

 From \eqref{Texplicit} and Lemma \ref{relationT}, we thus obtain an explicit, cancellation-free formula for the antipode $ S\big(\lx_{(a, 
 b)}\big)$, with coefficients given recursively in terms of Euler's totient function, as below. 

\begin{theorem}\label{mainantipodethm}
 For positive integers $a$ and $b$, let 
 $\mathcal{C}^{(a, b)}_{1} := -1 $ 
 and let 
\begin{equation}\label{displayfinalcoef}
 \mathcal{C}^{(a, b)}_{i} := -\left( 
 T(a+i,i) 
 + \sum_{j=1}^{i-1} \mathcal{C}^{(a, b)}_{j}
 T(a+i,i-j) \right) 
\end{equation} 
 for all possible indices $i$. Then 
\begin{equation*}
 S\big(\lx_{(a, b)}\big) = (-1)^{a+b} \big( E_{(b, a)} + \mathcal{C}_{1}^{(a, b)} E_{(b-1, a + 1)} + \mathcal{C}_{2}^{(a, b)} E_{(b-2, a + 2)}
 + \cdots + \mathcal{C}_{b}^{(a, b)} E_{(a+b)} \big). 
 \end{equation*} 
\end{theorem}

\begin{proof}
 This follows from Theorem \ref{mainccoeff} and \eqref{antipodeE} and Lemma \ref{relationT}. 
\end{proof}

\begin{example}
 We obtain the antipode evaluation $$ S\big( \lx_{33} \big) = E_{33} - E_{24} - E_{15}$$ using Theorem \ref{mainantipodethm}, and similarly 
 for $$ S\big( \lx_{42} \big) = E_{24} - E_{15} - 2 E_{6}. $$
\end{example}

 A similar approach, relative to our proof of Theorem \ref{mainantipodethm}, can be used to obtain cancellation-free formulas for the 
 antipodes of $\lx$-basis indexed by compositions with a fixed number of rows greater than $2$, but, as is the case with the immaculate 
 basis, this leads to more and more complicated formulas, and we encourage the exploration of such higher-order antipode formulas. 
 
 It would be desirable to apply Theorem \ref{mainantipodethm} to obtain a combinatorial interpretation for the coefficients of the form 
 shown in \eqref{displayfinalcoef}, and to apply Theorem \ref{mainantipodethm} to obtain a cancellation-free 
 expansion of $S(\lx_{(a, b)})$ in the $\lx$-basis. We leave these topics to a future study, and further research problems 
 concerning the $\lx$-basis are given in Section \ref{Conclusionsec} below. 
 
\section{Conclusion}\label{Conclusionsec}
 We conclude with open problems and topics for future research. To begin with, there are the natural questions regarding bases of 
 $\textsf{NSym}$ and $\textsf{QSym}$   and potential generalizations of results that apply to the immaculate and dual immaculate bases. Is 
 there a Jacobi--Trudi-like rule that the $\lx$-basis satisfies? How could the $\lx$-basis be constructed using analogues of the Bernstein 
 operators \cite{BergBergeronSaliolaSerranoZabrocki2014} used to construct the immaculate basis? How could a 
 Littlewood--Richardson-like product rule be determined for the $\lx$-basis? Letting $\chi\colon \textsf{NSym} \to \textsf{Sym}$ denote the usual projection morphism from $\textsf{NSym}$ and $\textsf{Sym}$, when 
 can the Schur-positivity of $\chi(\lx_{\alpha})$ be guaranteed? How can $\chi(\lx_{\alpha})$ be evaluated for a given 
 composition $\alpha$? 
 What is the expansion of a given Schur function into the dual lexical basis of $\textsf{QSym}$? Immaculate tableaux also relate closely to tabloids, so we may ask how the representation-theoretic applications of tabloids could be generalized using lexical tableaux, or other composition tableaux 
 analogues of tabloids. 
 
 Our definition of lexical tableaux and our construction of bases of $\textsf{QSym}$ and $\textsf{NSym}$ using such tableaux were based 
 on the analogue in Theorem \ref{theoremltstirling} of the summation identity in Theorem \ref{thm:snk} involving the number of standard 
 immaculate tableaux of a fixed height and of a fixed order. This summation identity was discovered experimentally using the OEIS, which 
 motivates further algebraic combinatorics-based explorations on summation identities involving expressions of the form $g^{\alpha}$. 
 In this direction, using the OEIS, we have experimentally discovered that \[ \sum_{\alpha \vDash
 n} g^{\alpha}\sum_{i} (-1)^{\alpha_i} = 
 \text{A000296}(n) = \sum_{k=0}^n (-1)^{n-k} \binom{n}{k}B(k) = B(n-1) - a(n-1), \] where $\text{A000296}$(n) counts number of set 
 partitions of $[n]$ without singletons, and that \[ \sum_{\alpha \vDash 
 n} g^{\alpha} \Bigg[ \sum_{i} \alpha_i (-1)^{\alpha_i+1} \Bigg] = 
 \text{A250105}(n) = n((-1)^{n-1} + \sum_{j=1}^{n-1} (-1)^{j-1}B(n-j-1)),\] for $n \geq 2$, where $\text{A250105}(n)$ counts the 
 number of partitions of $n$ with exactly one circular succession, and that \[ \sum_{\alpha \vDash 
 n} g^{\alpha} \ell(\alpha) 
 \sum_{i} \alpha_i (-1)^{\alpha_i +1} = \text{A052889}(n) = n B(n-1), \] where $\text{A052889}(n)$ is the number of rooted set partitions 
 of $[n]$. How could these experimentally discovered results be proved and applied by analogy with our results related to 
 Theorem \ref{theoremltstirling}? What similar relations might hold using standard lexical tableaux?
 
 \subsection*{Acknowledgements}
 The first author was supported through an NSERC Discovery Grant. 

\begin{thebibliography}{99}

\bibitem{AdinGesselReinerRoichman2021}%%%%%%%%%%%%%%%q
 \textsc{R.\ M.\ Adin, I.\ M.\ Gessel, V.\ Reiner, and Y.\ Roichman}, 
 Cyclic quasi-symmetric functions, 
 \emph{Israel J. Math.} {\bf 243(1)} (2021), 437--500. 

\bibitem{AguiarBergeronSottile2006}%%%%%%%%%%%%%%%
 \textsc{M.\ Aguiar, N.\ Bergeron, and F.\ Sottile}, 
 Combinatorial {H}opf algebras and generalized {D}ehn-{S}ommerville relations, 
 \emph{Compos.\ Math.} {\bf 142(1)} (2006), 1--30. 

\bibitem{AllenHallamMason2018}%%%%%%%%%%%%%%%
 \textsc{E.\ E.\ Allen, J.\ Hallam, and S.\ K.\ Mason}, 
 Dual immaculate quasisymmetric functions expand positively
 into {Y}oung quasisymmetric {S}chur functions, 
 \emph{J.\ Combin. Theory Ser. A} {\bf 157} (2018), 70--108. 

\bibitem{AllenMason2025}%%%%%%%%%%%%%%%
 \textsc{E.\ E.\ Allen and S.\ K.\ Mason}, 
 A combinatorial interpretation of the noncommutative inverse {K}ostka matrix, 
 \emph{Electron.\ J. Combin.} {\bf 32(1)} (2025), Paper No. 1.39, 48. 

 \bibitem{BenedettiSagan2017}%%%%%%%%%%%%%%%
 \textsc{C.\ Benedetti and B.\ E.\ Sagan}, 
 Antipodes and involutions, 
 \emph{J.\ Combin. Theory Ser. A} {\bf 148} (2017), 275--315. 

\bibitem{BergBergeronSaliolaSerranoZabrocki2015}%%%%%%%%%%%%%%%
 \textsc{C.\ Berg, N.\ Bergeron, F.\ Saliola, L.\ Serrano, and M.\ Zabrocki}, 
 Indecomposable modules for the dual immaculate basis of
 quasi-symmetric functions, 
 \emph{Proc. Amer. Math. Soc.} {\bf 143(3)} (2015), 991--1000. 

\bibitem{BergBergeronSaliolaSerranoZabrocki2014}%%%%%%%%%%%%%%%
 \textsc{C.\ Berg, N.\ Bergeron, F.\ Saliola, L.\ Serrano, and M.\ Zabrocki}, 
 A lift of the Schur and {H}all-{L}ittlewood bases to
 non-commutative symmetric functions, 
 \emph{Canad. J. Math.} {\bf 66(3)} (2014), 525--565. 

\bibitem{BergBergeronSaliolaSerranoZabrocki2017}%%%%%%%%%%%%%%%q
 \textsc{C.\ Berg, N.\ Bergeron, F.\ Saliola, L.\ Serrano, and M.\ Zabrocki}, 
 Multiplicative structures of the immaculate basis of non-commutative symmetric functions, 
 \emph{J.\ Combin.\ Theory Ser. A} {\bf 152} (2017), 10--44. 

 \bibitem{BergeronSanchezOrtegaZabrocki2016}%%%%%%%%%%%%%%%
 \textsc{N.\ Bergeron, J.\ S\'anchez-Ortega, and M.\ Zabrocki}, 
 The {P}ieri rule for dual immaculate quasi-symmetric functions, 
 \emph{Ann.\ Comb.} {\bf 20(2)} (2016), 283--300. 

\bibitem{Campbell2023}%%%%%%%%%%%%%%%q
 \textsc{J.\ M.\ Campbell}, 
 On antipodes of immaculate functions, 
 \emph{Ann.\ Comb.} {\bf 27(3)} (2023), 579--598. 

\bibitem{CassaigneFiciSciortinoZamboni2017}%%%%%%%%%%%%%%%
 \textsc{J.\ Cassaigne, G.\ Fici, M.\ Sciortino, and L.\ Q.\ Zamboni}, 
 Cyclic complexity of words, 
 \emph{J. Comb.\ Theory, Ser. A} {\bf 145} (2017), 36--56. 

 \bibitem{Daugherty2024}%%%%%%%%%%%%%%%
 \textsc{S.\ Daugherty}, 
 A generalization of the dual immaculate quasisymmetric functions in partially commutative variables, 
 \emph{Electron. J.\ Comb.} {\bf 31(2)} (2024), $\#$2.7, 59. 

\bibitem{DascalescuNastasescuRaianu2001}%%%%%%%%%%%%%%%
 \textsc{S.\ D\u asc\u alescu, C.\ N\u ast\u asescu, and \c S.\ Raianu}, 
 Hopf algebras, Marcel Dekker, Inc., New York (2001). 

\bibitem{GaoYang2016}%%%%%%%%%%%%%%%
 \textsc{A.\ L.\ L.\ Gao and A.\ L.\ B.\ Yang}, 
 A bijective proof of the hook-length formula for standard immaculate tableaux, 
 \emph{Proc.\ Amer.\ Math. Soc.} {\bf 144(3)} 
 (2016), 989--998. 

\bibitem{GelfandKrobLascouxLeclercRetakhThibon1995}%%%%%%%%%%%%%%%q
 \textsc{I.\ M.\ Gelfand, D.\ Krob, A.\ Lascoux, B.\ Leclerc, V.\ S.\ Retakh, and J.-Y.\ Thibon}, 
 Noncommutative symmetric functions, 
 \emph{Adv.\ Math.} {\bf 112(2)} 
 (1995), 218--348. 

\bibitem{Gessel1984}%%%%%%%%%%%%%%%
 \textsc{I.\ M.\ Gessel}, 
 Multipartite {$P$}-partitions and inner products of skew {S}chur functions, 
 \emph{Combinatorics and algebra ({B}oulder, {C}olo., 1983)}, 
 {\bf 34} Amer. Math. Soc., Providence, RI (1984), 289--317. 

\bibitem{gilbert1961symmetry}%%%%%%%%%%%%%%%
 \textsc{E.\ N.\ Gilbert and J. Riordan},
 Symmetry types of periodic sequences,
 \emph{Illinois J.\ Math.},
 {\bf 5(4)} (1961), 657--665.

\bibitem{Li2018}%%%%%%%%%%%%%%%q
 \textsc{S.\ X.\ Li}, 
 Structure constants for immaculate functions, 
 \emph{Ann.\ Comb.} {\bf 22(2)} (2018), 347--361. 

 \bibitem{LoehrNiese2021}%%%%%%%%%%%%%%%q
 \textsc{N.\ A.\ Loehr and E.\ Niese}, 
 Combinatorics of the immaculate inverse {K}ostka matrix, 
 \emph{Algebr.\ Comb.} {\bf 4(6)} (2021), 1119--1142. 

 \bibitem{Macdonald1995}%%%%%%%%%%%%%%%q
 \textsc{I.\ G.\ Macdonald}, 
 Symmetric functions and {H}all polynomials, 
 The Clarendon Press, Oxford University Press, New York (1995). 

 \bibitem{MasonSearles2021}%%%%%%%%%%%%%%%q
 \textsc{S.\ Mason and D.\ Searles}, 
 Lifting the dual immaculate functions, 
 \emph{J.\ Combin. Theory Ser. A} {\bf 184} (2021), Paper No. 105511, 52. 

\bibitem{MilnorMoore1965}%%%%%%%%%%%%%%%q
 \textsc{J.\ W.\ Milnor and J.\ C.\ Moore}, 
 On the structure of {H}opf algebras, 
 \emph{Ann.\ of Math. (2)} {\bf 81}
 (1965), 211--264. 

\bibitem{NieseSundaramvanWilligenburgWang2024}%%%%%%%%%%%%%%%
 \textsc{E.\ Niese, S.\ Sundaram, S.\ van Willigenburg, and S.\ Wang}, 
 Pieri rules for skew dual immaculate functions, 
 \emph{Can.\ Math.\ Bull.} {\bf 67(4)} (2024), 902--914. 

\bibitem{NovelliThibonToumazet2018}%%%%%%%%%%%%%%%q
 \textsc{J.-C.\ Novelli, J.-Y.\ Thibon, and F.\ Toumazet}, 
 Noncommutative {Bell} polynomials and the dual immaculate basis, 
 \emph{Algebr.\ Comb.} {\bf 1(5)} (2018), 653--676. 

\bibitem{Sagan2001}%%%%%%%%%%%%%%%
 \textsc{B.\ E.\ Sagan}, 
 The symmetric group, 
 Springer-Verlag, New York (2001). 

\bibitem{SunHu2018}%%%%%%%%%%%%%%%
 \textsc{B.\ Y.\ Sun and Y.\ Hu}, 
 A probabilistic method for the number of standard immaculate tableaux, 
 \emph{Rocky Mountain J.\ Math.} {\bf 48(6)} 
 (2018), 2087--2097. 

\end{thebibliography}

\end{document}
